\section{26 июня 1926 г.}

\lp{20260606_184237830}

Москва, 26 июня 1926 г.

Милый и дорогой Коля, письмо твое от 15 июня я получил 24-го, и хотя я и
предчувствовал, что ты не согласишься взять на себя обязанности профессора по
классу композиции, тем не менее отказ твой очень огорчил меня, Алекcандра
Борисовича, а также всех тех из профессоров консерватории, которые очень
надеялись видеть тебя в нашей среде. Доводы, тобой приведенные, далеко не все
убедительны (кроме последнего, а именно, что занятия эти отнимут у тебя время).
Первый довод о твоей, якобы, неподготовленности к этим занятиям основан на
неверном представлении о теперешнем ходе занятий на этом факультете.

\lp{20260606_184248627}

Если бы дело шло о «подвижном контрапункте» Танеева, то ты был бы прав, но ведь
все дисциплины, которые во времена Танеева являлись фундаментом и основой, в
настоящее время проходятся лишь вскользь, причем каждому профессору
предоставлено право вести дело по его усмотрению.

В этом имеются и хорошие и дурные стороны. Я лично глубоко уважаю и люблю
Сергея Ивановича, всегда боялся его, как педагога, и не решился, в свое время,
поступить к нему в класс; ты ушел от него, не пробыв и года и думаю, что
раскаяние тебя не тревожит, т.к. его подход к делу охлаждал учащихся, а иногда
вызывал и отвращение к этим занятиям.

\lp{20260606_184248627}

Всего правильнее, может быть, было бы закрыть совсем этот факультет, но так как
это невозможно, то надо найти какой-то разумный выход. Главное -- это выбор
руководителя. Можно многому научить и от многого предостеречь хотя бы только
советами, беседой и просмотром как работ учащихся, так и анализом настоящих
сочинений, и я уверен -- будь ты руководителем такого класса, польза для дела
будет огромная (вернее сказать «была бы», т.к. ты отказался).

Ну что же делать! Кандидатом теперь называют Жиляева, который в свое время
пресытившись подвижным контрапунктом, ищет новых путей,

\lp{20260606_184237830}

\noindent
и хотя сам совсем не сочиняет (это делает ему честь), но имеет и имел за
последние годы ряд учеников, успевших выдвинуться своей «левизной» и смелым
обращением с гармонией и прочими устарелостями.

Ты пишешь мне, что слышал о переходе многих учеников покойного Катуара еще при
жизни к другим преподавателям, и спрашиваешь, верно ли это? К сожалению, это
совершенно верно, и объясняется это явление очень просто. Егор Львович отдавал
все свои силы преподаванию; он к урокам готовился; тщательнейшим образом изучал
работы учащихся, и добивался того, что малоспособный ученик начинал в глубине
души сознавать

\lp{20260606_184300910}

\noindent
свою недостаточную одаренность и, не имея мужества бросить занятия музыкой,
переходил к другому преподавателю, утешая себя тем, что Катуар стар и не сумел
разгадать заложенных в нем талантов. А остальные профессора (кроме Г. Конюса)
смотрят на дело более спокойно, «трезво»: ни от каких эксцессов за голову не
хватаются, слегка умеряют пыл зарвавшихся не вмеру новаторов и смотрят так, что
чему быть, того не миновать, смотря на вещи глазами гоголевского Христиана
Ивановича из «Ревизора».

Учащимся это нравится, и они часто не осознавая, то пишут ерунду, продолжают
это делать.

\lp{20260606_184307524}

\noindent
вполне уверенные, что творят совершенно новую музыку.

Пантелеймон Иванович читал мне твое письмо, написанное как ответ Анат. Ник. на
его казуистический вопрос о том, что такое музыка. Так вот, содержание этого
твоего письма и могло бы лечь в основу программы 1-го курса творческого
факультета. Я чувствую, что нехорошо делаю, продолжая тебя волновать, т.к. ты
прекрасно сознаешь правоту и основания моего желания видеть именно тебя во
главе творческого факультета в Москве.
Ну довольно об этом.

Относительно семьи Катуара, дело устроилось в общем хорошо. Консерватория

\lp{20260606_184307524}

\noindent
получила официальное разрешение выдать семье жалование покойного Егора Львовича
за все летние месяцы. Таким образом его дети обеспечены до октября, а в октябре
А.В. Луначарский обещал пенсию. Если ты не писал еще Сергею Васильевичу по
этому поводу, то не пиши. Аня переехала в 2 комнаты своего отца, а свою комнату
отдала. Летом они живут на Сходне, так что в общем их положение сейчас не
плохо. Александр Борисович 8 июня давал концерт с Райским из сочинений Катуара.
Концерт этот, несмотря на африканскую погоду, дал все же некоторый доход,
который также переведен детям Егора Львовича. В концерте было исполнено:

\lp{20260606_184300910}

\noindent
фортепианное трио, фортепианный квартет и ряд романсов. Последним романсом был
романс на слова «смерть убаюкай меня», произведший большое впечатление как
музыкальное, так и исполнением желания покойного: ведь смерть действительно его
убаюкала; он уснул здоровым и больше не проснулся, причем лицо его и положение
тела свидетельствовали о том, что он не заметил как глубокий сон перешел в сон
вечный. Через несколько дней после смерти Егора Львовича скончался священник
Металлов, до конца жизни читавший в консерватории курс истории церковной
музыки. Странно, в Москве нет ни одного заместителя на эту вакансию, так что
курс этот будет заменен курсом палеографии, который будут читать
Иванов-Борецкий.

\lp{20260606_184505012}

На днях тебе будет писать, а может уже и написал письмо, Цейтлин, который
желает получить твое согласие на участие в концертах «Персимфанса» на ближайший
сезон. Имей ввиду, что «Росфил», с которым ты заключил договор, находится в
данное время в стадии реорганизации и, повидимому, в состав этого нового
учреждения войдут и новые лица. Думается мне, что ввиду этих новых
обстоятельств, ты не отказывайся от участия в «Персимфансе» до получения тобой
из нового «Росфила» подтверждения на заключенный тобой с Б. Б. Красиным договор.

\lp{20260606_184509749}

На тот случай, если бы ты захотел дать ряд концертов от себя (кроме концертов
«Росфила» и «Персимфанса» с твоим участием), то необходимо будет тебе летом же
до сентября озаботиться записью на Большой Зал консерватории, т.к. свободных
дней будет очень мало, а именно: понедельники -- «Персимфанс», вторники, среды,
четверга, субботы и воскресенья вечером -- Кино, пятницы и воскресенья утром --
концерты. Все эти вопросы надо решить теперь же. Если тебе понадобятся
какие-либо разъяснения, то я тебе отвечу тот час же, т.к., живя в
консерватории, встречаю постоянно Райского (можешь писать также и ему по этим
вопросам). Ведь он

\lp{20260606_184545580}

\noindent
является заведующим хозяйственной частью и все эти вопросы находятся всецело в
его компетенции.

Ты меня спрашивал в твоем письме к кому обратиться по вопросу о корректурах
твоих сочинений. Я думаю, что всего лучше к заведующему сектором А. Н. Юровскому
или к его секретарю Агнии Алексеевне Алявдиной.

Текущей весной на целом ряде ученических концертов и экзаменах я с наслаждением
прослушал большую часть твоих фортепианных произведений. Играли сонаты g-moll,
c-moll (сказки), c-moll (из цикла), все три сонаты из триады, соната
reminiscenza, ряд пьес из \rem{...}, соната e-moll(!) (у Фейнберга),

\lp{20260606_184548735}

концерт (очень хорошо), новеллы, скрипичная соната, импровизации, ряд сказок,
несколько раз похоронный марш и \rem{...} b-moll и h-moll. Ну словом, из всех
современных композиторов в ф.т. отделе консерватории, ты занял первое место и
вот мне и хотелось, чтобы на творческом факультете дело обстояло бы также, а ты
этого не хочешь.

Поблагодари Анну Михайловну за ее письмо. Очень жалею, что ты и Анна Михайловна
ничего не пишите о себе, а ведь вы же все-таки два раза пересекли Атлантический
океан (ты ни разу в своих письмах не писал о впечатлениях от переезда в Америку
и о самой Америке). Наконец, я всякий раз забываю тебя спросить: почему именно
ты избрал своим местом пребывания французскую деревню, почему не немецкую или
еще какую. Вы оба наверное по-французски теперь так наловчились говорить, что
при вас, пожалуй, и не решишься говорить на этом языке.

\lp{20260606_184526888}

Занятия в консерватории, слава Богу, наконец закончились (23 июня), и я начинаю
понемногу приходить в себя. Нынешний год я был очень занят в консерватории: у
меня был огромный камерный класс, ф.т. класс и класс органа, и если к этому
прибавить огромное количество ученических концертов, вечеров и невыносимых для
меня заседаний, то поневоле приходится удивляться выносливости человека и его
неумения жить по-человечески. Как бы то ни было, теперь у меня два месяца в
полном моем распоряжении.

Анна Михайловна спрашивала в письме как поживает Маня: она в общем такая же
осталась как и была при жизни мамаши. В 21-м году она вышла замуж

\lp{20260606_184533593}

Муж ее человек хороший, но с большими странностями. Если про Маню, всякий
знающий ее скажет, что психика ее неуравновешена и не совсем в порядке, то про
мужа ее приходится сказать то же самое, стой лишь разницей, что у Мани все
явления в жизни окрашиваются в пессимистические тона, а у мужа ее наоборот, но
также с явным уклоном к так называемой ненормальности. Вижу я их довольно
редко, но все-таки иногда захожу; заходит к нам и Маня. Ольгу с зятем вижу
довольно часто. Они, также как и мы, сидят в Москве. Паша захворал недавно
стоматитом (так называется воспаление

\lp{20260606_184533593}

\noindent
зева, рта и языка). Болезнь длительная и очень изнурительная. Соня с Андрюшей
уехали на дачу (последнее время Соня себя чувствует очень утомленной и слабой,
а Андрюша, хотя и выглядит всегда бодрым, но впечатление производит человека
подвинченного и очень нервного. В умственном отношении он развился чрезвычайно
и очень напоминает Милю в молодости. Паша наш остался без изменений (\rem{...}
только борода) и с тем же упорством звонит каждый день и рвет свою одежду. Дома
бывает редко. В отношении психики стал спокойнее; явления ненормальности
сократились значительно. Екатерина Петровна держится крепостью духа; детей
масса, отдыха нет да и жизни \rem{...}, в жару в городе очень тяжело.

\lp{20260606_184526888}

Ну словом, как говорится, пока живы и здоровы.

С нетерпением буду ждать твои скрипичные сонаты, которые надеюсь ты вышлешь до
твоего приезда в Москву. Передаю мой и Екатерины Петровны сердечный поклон Анне
Михайловне и поцелуй от меня крепко Милю. Скажи Льву Эдуардовичу и Ольге
Николаевне, что очень часто вспоминаю их, и если не пишу письма, то
исключительно из врожденной антипатии к этому занятию; кланяйся им от меня и
Екатерины Петровны.

Крепко тебя обнимаю и целую. Любящий тебя А. Гедике.
